% !TeX root = ../main.tex
% ----------------------------------------------------------------------
% Paquetes, colores, geometría y macros matemáticas
% ----------------------------------------------------------------------

% --- Fuentes y codificación (válido para XeLaTeX, LuaLaTeX y pdfLaTeX) ---
\usepackage{iftex}
\ifPDFTeX
  \usepackage[T1]{fontenc}
  \usepackage[utf8]{inputenc}
  \usepackage{lmodern}
\else
  \usepackage{fontspec}
\fi
\usepackage[spanish,es-nodecimaldot]{babel}
\spanishdeactivate{<>}
\AtBeginDocument{\shorthandoff{<>}}
\usepackage{microtype}

% --- Matemática ---
\usepackage{amsmath, amssymb, amsthm, mathtools}
\usepackage{mathrsfs}
\usepackage{cancel}
\usepackage{interval}

% --- Utilidades ---
\usepackage{graphicx}
\usepackage{enumitem}
\usepackage[table]{xcolor}
\usepackage{tikz}
\usepackage{pgfplots}
\pgfplotsset{compat=1.18}

% --- Colores (los mismos que en la portada del final de Temas) ---
\definecolor{darkblue}{RGB}{0, 51, 102}
\definecolor{accentblue}{RGB}{0, 102, 204}
\definecolor{mediumgray}{RGB}{100, 100, 100}
\definecolor{lightgray}{RGB}{240, 240, 240}

% --- Geometría y párrafos ---
\usepackage{geometry}
\geometry{
  a4paper,
  left=2.8cm, right=2.8cm,
  top=3.2cm, bottom=3.2cm,
  headheight=15pt
}
\setlength{\parindent}{0pt}
\setlength{\parskip}{0.6em plus 0.3em minus 0.2em}
\setlist{itemsep=0.2em, topsep=0.3em, parsep=0pt}

% --- Conjuntos numéricos y atajos ---
\newcommand{\N}{\mathbb{N}}
\newcommand{\Z}{\mathbb{Z}}
\newcommand{\Q}{\mathbb{Q}}
\newcommand{\R}{\mathbb{R}}
\newcommand{\C}{\mathbb{C}}
\newcommand{\K}{\mathbb{K}}
\newcommand{\bb}[1]{\mathbb{#1}}
\newcommand{\e}{\varepsilon}
\renewcommand{\O}{\varnothing}

% --- Operadores ---
\DeclareMathOperator{\card}{card}
\DeclareMathOperator{\id}{id}
\DeclareMathOperator{\Img}{Im}
\DeclareMathOperator{\diam}{diam}
\DeclareMathOperator{\dist}{dist}
\DeclareMathOperator{\interior}{int}
\newcommand{\cl}[1]{\overline{#1}}
\newcommand{\abs}[1]{\left\lvert #1 \right\rvert}
\newcommand{\norm}[1]{\left\lVert #1 \right\rVert}
\newcommand{\dd}{\mathop{}\!\mathrm{d}}

% --- Integrales superior e inferior ---
\def\upint{\mathchoice%
  {\mkern13mu\overline{\vphantom{\intop}\mkern7mu}\mkern-20mu}%
  {\mkern7mu\overline{\vphantom{\intop}\mkern7mu}\mkern-14mu}%
  {\mkern7mu\overline{\vphantom{\intop}\mkern7mu}\mkern-14mu}%
  {\mkern7mu\overline{\vphantom{\intop}\mkern7mu}\mkern-14mu}%
  \int}
\def\lowint{\mkern3mu\underline{\vphantom{\intop}\mkern7mu}\mkern-10mu\int}
